\section{Launch Exit Conditions and Orbital Insertion}
\label{sec:exit}

\subsection{Exit Velocity Definitions and Earth's Rotation}
A critical distinction must be made between platform-relative velocity and inertial (Earth-centered) velocity. The platform moves with Earth's rotation, providing an eastward velocity boost. At the Equator, Earth's rotational velocity is:
\begin{equation}
v_{rot} = \omega R_E \cos\phi
\end{equation}
where $\omega = 7.292 \times 10^{-5}$ rad/s is Earth's angular velocity, $R_E = 6371$ km is Earth's radius, and $\phi$ is latitude. At the Equator ($\phi = 0$):
\begin{equation}
v_{rot} = (7.292 \times 10^{-5}) \times 6371000 \times \cos(0) \approx 465  \text{m/s}
\end{equation}
This represents a \textbf{significant free velocity boost} that reduces the required platform-relative acceleration.

If the launcher fires eastward (prograde), the inertial exit velocity becomes:
\begin{equation}
v_{inertial} = v_{platform-relative} + v_{rot}
\end{equation}
For a platform-relative speed of 7.3 km/s eastward, the inertial velocity would be approximately 7.765 km/s. However, if firing westward (retrograde), the inertial velocity would be $7.3 - 0.465 = 6.835$ km/s, demonstrating the importance of launch azimuth.

\subsection{Orbital Mechanics and Required Delta-v}
Reaching 7.3 km/s platform-relative speed (or 7.765 km/s inertial with eastward launch) does \emph{not} guarantee a stable orbit. Orbital insertion requires sufficient specific orbital energy to achieve a closed trajectory that avoids immediate reentry.

The specific orbital energy $\epsilon$ is given by:
\begin{equation}
\epsilon = \frac{v^2}{2} - \frac{\mu}{r}
\end{equation}
where $v$ is inertial speed, $\mu = 3.986 \times 10^{14}$ m$^3$/s$^2$ is Earth's gravitational parameter, and $r$ is distance from Earth's center.

For a circular orbit at altitude $h$, the required orbital velocity is:
\begin{equation}
v_{circ} = \sqrt{\frac{\mu}{R_E + h}}
\end{equation}
At 400 km altitude (a common LEO destination):
\begin{equation}
v_{circ} = \sqrt{\frac{3.986 \times 10^{14}}{6371000 + 400000}} \approx 7669  \text{m/s} \approx 7.67  \text{km/s}
\end{equation}
At 45 km altitude (launch altitude), the circular orbital velocity would be:
\begin{equation}
v_{circ,45km} = \sqrt{\frac{3.986 \times 10^{14}}{6371000 + 45000}} \approx 7904  \text{m/s} \approx 7.90  \text{km/s}
\end{equation}

However, launching horizontally at 45 km altitude does not place the payload on a circular orbit at that altitude because:
\begin{itemize}
    \item \textbf{Perigee at launch altitude}: The trajectory's perigee (lowest point) would be at or near the launch altitude
    \item \item \textbf{Apogee requires energy}: To raise apogee to a useful altitude (e.g., 400-800 km for LEO constellations), additional $\Delta v$ is required at apogee
\end{itemize}

Using the vis-viva equation for an elliptical orbit:
\begin{equation}
v = \sqrt{\mu \left( \frac{2}{r} - \frac{1}{a} \right)}
\end{equation}
where $a$ is the semi-major axis.

For a transfer orbit with perigee at 45 km altitude and apogee at 400 km altitude:
\begin{equation}
r_p = 6371000 + 45000 = 6416000  \text{m}
\end{equation}
\begin{equation}
r_a = 6371000 + 400000 = 6771000  \text{m}
\end{equation}
\begin{equation}
a = \frac{r_p + r_a}{2} = \frac{6416000 + 6771000}{2} = 6593500  \text{m}
\end{equation}
The velocity at perigee (launch point) needed for this transfer orbit is:
\begin{equation}
v_{perigee} = \sqrt{\mu \left( \frac{2}{r_p} - \frac{1}{a} \right)} = \sqrt{3.986 \times 10^{14} \left( \frac{2}{6416000} - \frac{1}{6593500} \right)} \approx 7968  \text{m/s}
\end{equation}
Compared to the circular orbital velocity at 45 km altitude (7.904 km/s), this represents only a small $\Delta v$ increase (~64 m/s) to achieve the transfer orbit. However, this analysis assumes the launch velocity is purely tangential and at the correct altitude.

More significantly, after reaching apogee at 400 km altitude, the velocity would be:
\begin{equation}
v_{apogee} = \sqrt{\mu \left( \frac{2}{r_a} - \frac{1}{a} \right)} \approx 6778  \text{m/s}
\end{equation}
To circularize at 400 km altitude requires an additional $\Delta v$ of:
\begin{equation}
\Delta v_{circ} = v_{circ,400km} - v_{apogee} \approx 7669 - 6778 = 891  \text{m/s}
\end{equation}

This demonstrates that even with a 7.3 km/s platform-relative speed (7.765 km/s inertial eastward), substantial additional propulsion is required for stable orbital insertion due to:
\begin{itemize}
    \item \textbf{Gravity losses}: Though reduced by horizontal launch, vertical climb still incurs losses
    \item \item \textbf{Atmospheric drag}: During ascent through remaining atmosphere
    \item \item \textbf{Inclination changes}: If launch azimuth doesn't match target orbital plane
    \item \item \textbf{Perigee raising}: To avoid rapid reentry from perigee at launch altitude
\end{itemize}

The total $\Delta v$ required from launch to stable LEO typically exceeds 500-1000 m/s beyond the orbital velocity at launch altitude, meaning the payload needs either:
\begin{itemize}
    \item \textbf{Onboard propulsion}: A kick stage or integrated propulsion system
    \item \item \item \textbf{Upper stage}: A separate propulsion module that separates after burnout
    \item \item \item \item \textbf{Rendezvous with existing infrastructure}: Such as an orbital transfer vehicle (OTV) or space tug
\end{itemize}

\subsection{Aerodynamic Heating and Loads at Tube Exit}
Upon exiting the accelerator tube, the payload encounters the stratospheric atmosphere (~0.3-0.01 kg/m$^3$ density at 40-50 km) at hypersonic speeds (Mach 20-25). This creates severe aerothermal and mechanical loads:

\subsubsection{Stagnation Temperature}
The theoretical stagnation temperature is given by:
\begin{equation}
T_0 = T_\infty \left(1 + \frac{\gamma-1}{2} M^2 \right)
\end{equation}
where $T_\infty$ is ambient temperature (~250 K), $\gamma = 1.4$ for air, and $M$ is Mach number.

At Mach 24 ($v = 7.3$ km/s, $a \approx 300$ m/s at 45 km):
\begin{equation}
T_0 = 250 \left(1 + \frac{0.4}{2} \times 24^2 \right) = 250 \left(1 + 0.2 \times 576 \right) = 250 \times 116.2 \approx 29,050  \text{K}
\end{equation}
This represents the \textit{theoretical} stagnation temperature assuming adiabatic compression. In reality, chemical dissociation, ionization, and radiation reduce the actual surface temperature, but values still exceed 2000 K, requiring significant thermal protection.

\subsubsection{Heat Flux}
Convective heat flux can be estimated using the Fay-Riddell equation:
\begin{equation}
\dot{q} \propto \sqrt{\frac{\rho}{R_n}} v^3
\end{equation}
where $\rho$ is atmospheric density, $R_n$ is nose radius, and $v$ is velocity.

Even with a blunt nose ($R_n \approx 0.1$ m), heat fluxes of 10-100 MW/m$^2$ are possible at these conditions, necessitating active cooling or ablative thermal protection systems.

\subsubsection{Static Pressure and Drag}
The dynamic pressure is:
\begin{equation}
q = \frac{1}{2} \rho v^2
\end{equation}
At $\rho = 0.1$ kg/m$^3$ and $v = 7300$ m/s:
\begin{equation}
q = \frac{1}{2} \times 0.1 \times (7300)^2 \approx 26.6  \text{MPa}
\end{equation}
This represents an enormous structural load on the payload's nose and leading edges.

The drag force is:
\begin{equation}
F_D = q C_D A
\end{equation}
For a payload with reference area $A = 0.1$ m$^2$ and drag coefficient $C_D \approx 0.5$:
\begin{equation}
F_D = 26.6 \times 10^6 \times 0.5 \times 0.1 \approx 1.33  \text{MN}
\end{equation
This drag force is comparable to the thrust force used for acceleration, meaning significant deceleration begins immediately after exit.

\subsubsection{Shock-Wave Interactions**
The payload will experience complex shock-wave interactions including:
\begin{itemize}
    \item \textbf{Bow shock}: Formed in front of the payload
    \item \item \item \textbf{Shock-boundary layer interaction}: Potential separation and heating spikes
    \item \item \item \item \item \textbf{Interference shocks**: From protrusions, windows, or surface irregularities
\end{itemize}
These can create localized heating and pressure peaks far exceeding average values.

\subsection{Service Module Architecture and Propulsion**
To achieve orbital insertion, the payload likely requires a service module providing post-launch propulsion. We analyze the implications of 1,000 g acceleration on this module and the overall spacecraft.

\subsubsection{Propulsion System Requirements**
Assuming a total $\Delta v$ requirement of 900 m/s for orbital insertion (gravity losses, drag, perigee raise), and using a propulsion system with specific impulse $I_{sp} = 300$ s (typical for monopropellant or low-performance bipropellant):
\begin{equation}
\Delta v = I_{sp} g_0 \ln\left(\frac{m_0}{m_f}\right)
\end{equation}
where $g_0 = 9.81$ m/s$^2$, $m_0$ is initial mass, and $m_f$ is final mass.

Solving for propellant mass fraction:
\begin{equation}
\frac{m_p}{m_0} = 1 - e^{-\Delta v / (I_{sp} g_0)} = 1 - e^{-900 / (300 \times 9.81)} \approx 1 - e^{-0.306} \approx 0.264
\end{equation}
Thus, approximately 26.4\% of the service module's initial mass must be propellant. For a 150 kg payload, if the service module masses 50 kg, it would require approximately 13.2 kg of propellant.

More advanced propulsion (e.g., $I_{sp} = 350$ s) reduces this to ~23.1\%. Electric propulsion ($I_{sp} = 1000-3000$ s) would require only a few percent propellant but introduces power generation and thermal management challenges.

\subsubsection{Guidance, Attitude Control, and Thermal Protection**
The service module must provide:
\begin{itemize}
    \item \textbf{Attitude control}: To maintain correct orientation during thrust vectoring
    \item \item \item \textbf{Guidance}: Navigation updates to correct for wind drift, atmospheric variations, and injection errors
    \item \item \item \item \textbf{Thermal protection}: Shielding from aerodynamic heating during ascent
    \item \item \item \item \item \textbf{Thrust vectoring**: Gimbaling or alternative methods to direct thrust
\end{itemize}
These systems add mass, complexity, and potential failure points.

\subsubsection{Perigee Raising and Orbital Lifetime**
Even if the payload reaches apogee at a useful altitude, if perigee remains at or near launch altitude (~45 km), atmospheric drag will cause rapid reentry (within hours to days). Raising perigee to at least 200-300 km is typically required for meaningful orbital lifetime (weeks to years), which may require additional $\Delta v$ beyond basic circularization.

\subsection{Payload Architecture Under 1,000 g Acceleration**
The 1,000 g acceleration environment imposes extreme demands on payload components and structure:

\subsubsection{Electronics**
Standard commercial electronics cannot survive 1,000 g. Requirements include:
\begin{itemize}
    \item \textbf{Component selection}: Military-spec or space-grade components rated for high g-loads
    \item \item \item \textbf{Potting and encapsulation**: Embedding components in epoxy or silicone to prevent movement
    \item \item \item \item \textbf{Board stiffeners**: Mechanical supports to prevent flexing and cracking
    \item \item \item \item \item \textbf{Minimal mass**: Reducing inertial forces ($F = ma$)
    \item \item \item \item \item \item \textbf{Reduced complexity**: Fewer parts and solder joints to minimize failure points
\endcite{hoffmann2010}
\end{itemize}
Testing at centrifuge facilities is required to validate survival.

\subsubsection{Batteries and Energy Storage**
Batteries face unique challenges under high g:
\begin{itemize}
    \item \textbf{Electrode delamination**: Separation of anode/cathode/separator layers
    \item \item \item \textbf{Short circuit risk**: From internal component movement
    \item \item \item \item \textbf{Temperature sensitivity**: Performance changes under load
    \item \item \item \item \item \textbf{Containment**: Robust casing to prevent rupture and leakage
\end{itemize}
Specialized high-g batteries or alternative energy storage (e.g., capacitors for short bursts) may be required.

\subsubsection{CubeSat Structures and Deployables**
Standard CubeSat structures (aluminum frames, PC/104 stacking) are not designed for 1,000 g. Considerations include:
\begin{itemize}
    \item \textbf{Monocoque designs**: Integrated structures that distribute load more effectively
    \item \item \item \textbf{Carbon fiber composites**: High specific strength and tailored stiffness
    \item \item \item \item \item \textbf{Internal bracing**: Triangulated structures to prevent buckling
    \item \item \item \item \item \item \textbf{Deployable systems**: Antennas, solar panels, and instruments must survive acceleration and deploy reliably afterward
\end{itemize}
Solar panels, in particular, face challenges with cell cracking and connection failures under high g.

\subsubsection{Modular Spacecraft ("Space Legos")**
To facilitate assembly and repair under constraints, a modular approach may be beneficial:
\begin{itemize}
    \item \textbf{Standardized interfaces**: Mechanical, electrical, and data connectors designed for high g
    \item \item \item \textbf{Autonomous docking**: Vision-based or sensor-guided mating without human intervention
    \item \item \item \item \item \textbf{Robotic joining**: Small manipulators or self-aligning features for micro-assembly
    \item \item \item \item \item \item \textbf{Electrical/data interfaces**: Reliable connectors that maintain integrity under vibration and thermal cycling
    \item \item \item \item \item \item \item \textbf{In-space manufacturing**: Additive manufacturing or assembly techniques for on-orbit repairs or upgrades
\end{itemize}
However, these systems add mass, complexity, and potential failure points, and their ability to function after 1,000 g exposure requires validation.

\subsubsection{Functional Integration and Testing**
Validating payload survivability requires:
\begin{itemize}
    \item \textbf{Component-level testing**: Centrifuge testing of individual parts
    \item \item \item \textbf{Subsystem validation**: Testing assembled boards and mechanisms
    \item \item \item \item \item \textbf{Integrated system testing**: Full payload vibration and shock profiles
    \item \item \item \item \item \item \textbf{Environmental validation**: Combined thermal, vacuum, and vibration profiles
    \item \item \item \item \item \item \item \textbf{Flight heritage**: Incremental testing through suborbital demonstrations
\endcite{NASA2011}
\end{itemize}
Without such validation, assumptions about payload survival remain speculative.

\subsection{Connections to Broader Spacecraft Design**
The 1,000 g environment shares similarities with and differences from other extreme acceleration environments:

\subsubsection{Comparison with Artillery and Gun-Launch Systems**
Historical systems like Project HARP achieved ~1,000 g for suborbital launches. Lessons include:
\begin{itemize}
    \item \textbf{Short duration**: Milliseconds-scale events allow different material responses than sustained loading
    \item \item \item \textbf{Axis alignment**: Critical for preventing wobble and transverse loads
    \item \item \item \item \item \textbf{Recovery systems**: Parachutes or other methods to survive impact with water or ground
\end{itemize}
However, orbital insertion requires sustained velocity management after launch, differing from pure ballistic trajectories.

\subsubsection{Comparison with Aircraft Carrier Arresting and Ejection Seats**
These systems experience high g-forces but for very short durations and with different orientation profiles. Lessons focus on:
\begin{itemize}
    \item \textbf{Human tolerance limits**: Relevant if considering crewed variants (though unlikely for nanosatellite launch)
    \item \item \item \item \textbf{Energy absorption**: Techniques for dissipating kinetic energy reliably
    \item \item \item \item \item \item \textbf{Guided tracking**: Maintaining orientation during high-g transients
\end{itemize}
The horizontal, sustained nature of the pneumatic accelerator presents distinct challenges.

\subsubsection{Implications for Payload Mass and Complexity**
The need to survive 1,000 g likely increases payload mass and decreases complexity:
\begin{itemize}
    \item \textbf{Structural mass fraction**: Increases due to thicker walls, stiffer materials, and internal bracing
    \item \item \item \textbf{Reduced payload fraction**: Less mass available for instruments, power, and communications
    \item \item \item \item \item \textbf{Increased reliability focus**: Emphasis on simplicity, redundancy, and failure tolerance
    \item \item \item \item \item \item \textbf{Testing overhead**: Significant investment in validation facilities and procedures
\end{itemize}
This tension between survivability and mission capability represents a key design trade-off.

\subsection{Research Questions and Open Challenges**
This section identifies critical open questions requiring investigation:

\subsubsection{Integrated Trajectory Analysis**
\begin{itemize}
    \item \textbf{What is the actual $\Delta v$ budget** from launch to stable orbit, considering gravity losses, drag, and inclination changes?
    \item \item \item \textbf{How do stratospheric winds and atmospheric variability** affect launch accuracy and required corrections?
    \item \item \item \item \item \textbf{What lift-to-drag ratio** is achievable or required during ascent for lift-assisted trajectories?
\end{itemize}

\subsubsection{Propulsion System Selection**
\begin{itemize}
    \item \textbf{What is the optimal trade-off** between specific impulse, thrust-to-weight ratio, storability, and reliability for the kick stage?
    \item \item \item \textbf{Can electric propulsion** be made viable despite power and thermal constraints?
    \item \item \item \item \item \textbf{How much propellant mass** is realistically needed, and what tankage fraction is acceptable?
\end{itemize}

\subsubsection{Thermal Protection System (TPS)**
\begin{itemize}
    \item \textbf{What materials** can survive the combined thermal, mechanical, and chemical environment** of hypersonic exit?
    \item \item \item \textbf{How much TPS mass** is required, and can it be integrated structure rather than added mass?
    \item \item \item \item \item \textbf{Can active cooling** be made lightweight and reliable enough for this application?
\end{itemize}

\subsubsection{Attitude Determination and Control System (ADCS)**
\begin{itemize}
    \item \textbf{What sensors** (star trackers, sun sensors, IMUs, GPS) can function** during and after the high-g launch transient?
    \item \item \item \item \textbf{How much control authority** is needed to counteract aerodynamic torques during ascent?
    \item \item \item \item \item \item \textbf{Can reaction wheels or control moment gyros** survive the launch transient, or are thrusters or magnetorquers required?
\end{itemize}

\subsubsection{Post-Launch Payload Functionality**
\begin{itemize}
    \item \textbf{What fraction of payload mass** must be dedicated to survivability** rather than mission functionality?
    \item \item \item \item \item \textbf{How do standard CubeSat buses** need to be modified** for high-g compatibility?
    \item \item \item \item \item \item \textbf{What testing protocols** are sufficient to qualify** a payload for 1,000 g launch?
\end{itemize}

\subsubsection{Modular and Reconfigurable Architectures**
\begin{itemize}
    \item \textbf{Can standardized interfaces** survive** the launch transient and enable** on-orbit assembly or repair?
    \item \item \item \item \item \textbf{What is the mass penalty** of modularity versus integrated design?
    \item \item \item \item \item \item \textbf{How does autonomous docking reliability** compare to pre-integrated systems under post-launch constraints?
\end{itemize}

This section has established that reaching 7.3 km/s does not guarantee orbital insertion, quantified the additional propulsion requirements, analyzed the severe exit environment, and detailed the implications of 1,000 g on payload architecture. The following section will examine recovery systems, economics, and conclude with a comprehensive feasibility assessment.
